% !TeX root = Report.tex
\chapter{Использование \texorpdfstring{\texttt{CSSA row.row}}{CSSA row.row} для подавления шума в случае одной прямой}

В данной главе рассматривается внутренний механизм работы варианта \texttt{CSSA row.row} в задаче выделения одной прямой. Основная цель состоит в том, чтобы описать, как данный метод подавляет шум, какие ограничения связаны с числом компонент и расположением прямой на изображении, а также почему компоненты с наибольшими сингулярными значениями иногда не имеют отношения к сигналу.

\section{Шумоподавление методом \texorpdfstring{\texttt{CSSA row.row}}{CSSA row.row}}

Во всех последующих иллюстрациях рассматривается дискретное изображение размера
$100\times100$. Пиксели изображения индексируются парами
\[
(x,y)\in \{1,\dots,100\}\times\{1,\dots,100\}.
\]
На этой сетке задается одна прямая вида
\[
y=2x-1.
\]
Так как изображение является дискретным, фактически рассматривается множество пикселей
\[
\{(x,2x-1)\colon x\in\{1,\dots,50\}\},
\]
то есть $x$ пробегает целые значения, для которых точка остается внутри изображения.

К изображению добавлен гауссовский белый шум со стандартным отклонением $\sigma=0.2$; дополнительно рассматривается более слабый шум с $\sigma=0.05$. Предобработка строится по схеме, указанной в главе~1: к каждой строке изображения применяется дискретное преобразование Фурье, затем в каждой строке выполняется реконструкция по одной компоненте \texttt{CSSA}, после чего применяется обратное преобразование Фурье.

На рис.~\ref{fig:row-row-basic} приведен типичный результат шумоподавления методом \texttt{CSSA row.row} для одной прямой. По умолчанию предполагается, что число оцениваемых прямых равно \(1\), а длина окна выбирается равной \(L = N/2\), где $N$ --- длина строки в матрице изображения.

\begin{figure}[H]
	\centering
	\includegraphics[width=0.84\textwidth]{cssa_row_row_denoising.pdf}
	\caption{Базовый пример шумоподавления методом \texttt{CSSA row.row} для одной прямой при $\sigma=0.2$}
	\label{fig:row-row-basic}
\end{figure}

Основная форма прямой при такой предобработке сохраняется, а значительная часть шума подавляется. Вместе с тем на изображении остаются слабые следы и отдельные смещенные точки. Для понимания природы этих эффектов необходимо рассмотреть ограничения метода более подробно.

\section{Ограничения по числу компонент и расположению прямой}

Пусть $p$-я строка изображения содержит один пиксель сигнала в столбце $j_p$ с интенсивностью $c_p>0$. Тогда соответствующий одномерный ряд имеет вид
\[
\mathbf{x}^{(p)} = c_p \mathbf{e}_{j_p} + \boldsymbol{\varepsilon}^{(p)},
\]
где $\mathbf{e}_{j_p}$ --- стандартный базисный вектор, а $\boldsymbol{\varepsilon}^{(p)}$ --- шумовая компонента. Для чистого сигнала после дискретного преобразования Фурье получаем
\[
\mathcal{F}\bigl(\mathbf{x}^{(p)}\bigr)
=
c_p \boldsymbol{\omega}_{j_p},
\qquad
\bigl(\boldsymbol{\omega}_{j_p}\bigr)_k
=
\exp\!\left(
-\frac{2\pi i}{n}(k-1)(j_p-1)
\right).
\]
Таким образом, чистая строка в частотном представлении задается одной комплексной экспонентой \cite{trickett2003,zhornikova2016}. Именно с этим связан выбор одной компоненты в дальнейшей реконструкции: для такой экспоненты траекторная матрица имеет ранг $1$.

Следующее утверждение дает точную формулировку этого факта.

\begin{statement}
Пусть задан комплексный ряд
\[
x_k = e^{\alpha(k-1)} e^{i\omega(k-1)}, \qquad k=1,2,\dots,N,
\]
где $\alpha,\omega \in \mathbb{R}$, и зафиксирована длина окна $L$,
удовлетворяющая условию $1<L<N$. Рассмотрим его траекторную матрицу
\[
\mathbf{X}=
\begin{pmatrix}
x_1 & x_2 & \dots & x_K\\
x_2 & x_3 & \dots & x_{K+1}\\
\vdots & \vdots & \ddots & \vdots\\
x_L & x_{L+1} & \dots & x_N
\end{pmatrix},
\qquad K=N-L+1.
\]
Тогда
\[
\operatorname{rank}(\mathbf{X})=1.
\]
\end{statement}

\begin{proof} Обозначим
\[
c=e^{\alpha+i\omega}.
\]
Элемент матрицы $\mathbf{X}$ с индексами $(i,j)$ равен
\[
X_{ij}=x_{i+j-1}=c^{i+j-2} = c^{i-1} c^{j-1}.
\]
Следовательно,
\[
\mathbf{X}=uv^{\top},
\]
где
\[
u=
\begin{pmatrix}
1\\ c\\ c^2\\ \vdots\\ c^{L-1}
\end{pmatrix},
\qquad
v=
\begin{pmatrix}
1\\ c\\ c^2\\ \vdots\\ c^{K-1}
\end{pmatrix}.
\]
То есть $\mathbf{X}$ представима в виде внешнего произведения двух векторов. Поэтому
\[
\operatorname{rank}(\mathbf{X}) \le 1.
\]
Поскольку оба вектора $u$ и $v$ ненулевые, имеем $\mathbf{X}\neq 0$. Следовательно,
\[
\operatorname{rank}(\mathbf{X})\neq 0.
\]
Следовательно,
\(
\operatorname{rank}(\mathbf{X})=1.
\)
\end{proof}

Для рассматриваемой задачи это утверждение применяется следующим образом. Если $\mathbf{e}_j$ --- единичный орт длины $N$, то после дискретного преобразования Фурье получаем ряд
\[
\bigl(\mathcal{F}(\mathbf{e}_j)\bigr)_k
=
\exp\!\left(
-\frac{2\pi i}{N}(j-1)(k-1)
\right),
\qquad k=1,\dots,N.
\]
Следовательно, условие теоремы выполняется при $\alpha=0$ и
\[
\omega = -\frac{2\pi}{N}(j-1).
\]
Поэтому траекторная матрица такого ряда имеет ранг $1$. Если яркость пикселя равна $c_k < 1$, то, в силу линейности преобразования Фурье,  $\bigl(\mathcal{F}(c_ke_j)\bigr) = c_k\bigl(\mathcal{F}(e_j)\bigr)$, поэтому ранг соответствующей траекторной матрицы не изменится.

Рассмотрим теперь, как ведет себя шумовая часть после перехода в частотную область. Пусть
\[
\boldsymbol{\varepsilon}\sim \mathcal{N}(0,\sigma^2 I_N),
\]
а $F$ --- матрица дискретного преобразования Фурье без дополнительной нормировки:
\[
F_{kl}
=
\exp\!\left(
-\frac{2\pi i}{N}(k-1)(l-1)
\right),
\qquad k,l=1,\dots,N.
\]
Тогда для строки с одним пикселем сигнала
\[
\mathbf{x}=c\mathbf{e}_j+\boldsymbol{\varepsilon}
\]
после ДПФ получаем
\[
F\mathbf{x}
=
cF\mathbf{e}_j + F\boldsymbol{\varepsilon}
=
c\boldsymbol{\omega}_j+\widetilde{\boldsymbol{\varepsilon}}.
\]
Вектор $\widetilde{\boldsymbol{\varepsilon}}=F\boldsymbol{\varepsilon}$ остается гауссовским, поскольку любая его линейная комбинация имеет вид
\[
a^{*}\widetilde{\boldsymbol{\varepsilon}}
=
a^{*}F\boldsymbol{\varepsilon}
=
(F^{*}a)^{*}\boldsymbol{\varepsilon},
\]
то есть является линейной комбинацией координат гауссовского вектора. При этом эрмитова ковариационная матрица равна
\[
\operatorname{Cov}(\widetilde{\boldsymbol{\varepsilon}})
=
\operatorname{Cov}(F\boldsymbol{\varepsilon})
=
\sigma^2 FF^{*}.
\]
Так как
\[
(FF^{*})_{kl}
=
\sum_{s=0}^{N-1}
\exp\!\left(
-\frac{2\pi i}{N}s(k-l)
\right)
=
\begin{cases}
N, & k=l,\\
0, & k\ne l,
\end{cases}
\]
то
\[
\operatorname{Cov}(F\boldsymbol{\varepsilon})
=
N\sigma^2 I_N.
\]
Следовательно, ДПФ не уничтожает шум, а переводит его в комплексную гауссовскую добавку с увеличенной в $N$ раз дисперсией в принятой ненормированной записи преобразования. Именно эта добавка нарушает точное свойство ранга $1$: траекторная матрица раскладывается в сумму сигнальной части ранга $1$ и шумовой части, которая в общем случае уже не имеет малого ранга.

Тем не менее в качестве приближения по-прежнему можно использовать одну компоненту. Обычно при этом выбирается компонента, соответствующая наибольшему сингулярному числу траекторной матрицы, тогда как компоненты шума связаны с меньшими сингулярными числами. Это наблюдение иллюстрируется на простом опыте с одной строкой длины $n=100$, содержащей один пиксель сигнала, для двух уровней гауссовского шума: $\sigma=0.2$ и $\sigma=0.05$.

\begin{figure}[H]
	\centering
	\begin{subfigure}{0.49\textwidth}
		\centering
		\includegraphics[width=\textwidth]{singular-values-02.pdf}
		\caption{$\sigma=0.2$}
	\end{subfigure}
	\hfill
	\begin{subfigure}{0.49\textwidth}
		\centering
		\includegraphics[width=\textwidth]{singular-values-005.pdf}
		\caption{$\sigma=0.05$}
	\end{subfigure}
	\caption{Первые двенадцать сингулярных значений для чистой и зашумленной строки длины $n=100$ после дискретного преобразования Фурье}
	\label{fig:one-row-singular-values}
\end{figure}

Для чистой строки ненулевой остается только первая компонента, тогда как шум делает ненулевыми и последующие сингулярные значения. При уровне шума $\sigma=0.2$ сингулярный спектр заметно растягивается. Это значит, что первые несколько компонент сингулярного разложения приносят вклад. Так как в реконструкции сохраняется только одна компонента с наибольшим сингулярным значением, при наличии шума оценка может смещаться относительно истинного сигнала. При меньшем шуме $\sigma=0.05$ доминирование первой компоненты выражено существенно сильнее: остальные сингулярные значения оказываются малы по сравнению с первым. В этом случае разложение даст более устойчивый результат.

Указанная однокомпонентная модель является корректной лишь тогда, когда в рассматриваемой строке действительно присутствует один элемент сигнала. В более общем случае строка может содержать несколько пикселей сигнала:
\[
\mathbf{x}^{(p)} = \sum_{\ell=1}^{q_p} c_{p\ell}\mathbf{e}_{j_{p\ell}} + \boldsymbol{\varepsilon}^{(p)},
\]
где $q_p$ --- число пикселей сигнала в $p$-й строке. После преобразования Фурье такая строка переходит в сумму $q_p$ комплексных экспонент. Поэтому реконструкция по числу компонент, меньшему чем $q_p$, может привести к потере части сигнала. Выбор несколько большего числа компонент менее опасен: сигнал при этом, как правило, сохраняется, хотя в реконструкцию попадает и дополнительный шум.

С содержательной точки зрения здесь следует различать две разные ситуации.
\begin{enumerate}
	\item Если на изображении присутствуют две прямые, а реконструкция выполняется так, как если бы прямая была одна, то одна из двух прямых может вообще не восстановиться. Это не является ограничением метода как такового, а указывает на неправильный выбор числа восстанавливаемых компонент.
	\item Если прямая почти горизонтальна, то из-за дискретизации некоторые строки содержат несколько соседних пикселей сигнала, а некоторые строки не содержат сигнала вовсе и состоят только из шума. Именно это уже является содержательным ограничением построчной схемы \texttt{CSSA row.row}, поскольку проблема возникает из-за самого расположения прямой на дискретной сетке.
\end{enumerate}

Следовательно, выбирать число компонент меньше истинного числа структур сигнала недопустимо, поскольку это приводит к неправильному восстановлению сигнала. Напротив, выбор несколько большего числа компонент возможен, хотя и приводит к возврату части шума. Отдельно от этого сохраняется ограничение, связанное с дискретизацией почти горизонтальных прямых; именно оно далее проявляется в эффектах растекания и пропуска пикселей сигнала.

Это видно и в однострочном эксперименте при сравнении реконструкций по группам $1$, $1{:}2$, $1{:}3$ и $1{:}4$.

\begin{figure}[H]
	\centering
	\includegraphics[width=0.88\textwidth]{one_row_reconstructed_rows.pdf}
	\caption{Реконструкция одной строки при сохранении одной, двух, трех и четырех компонент}
	\label{fig:one-row-reconstructed}
\end{figure}

Реконструкция по одной компоненте дает главный максимум в правильной позиции. При добавлении второй, третьей и четвертой компонент в строке появляются новые значения, удаленные от истинного пикселя сигнала. Для одной прямой увеличение числа компонент поэтому не следует рассматривать как основной способ улучшения результата.

Даже при реконструкции по одной компоненте остаются два характерных дефекта, о которых было сказано выше: растекание и смещение.

Первый дефект не меняет положения основного максимума, однако расширяет множество положительных значений в восстановленной строке. Второй дефект существенно опаснее, поскольку уже на уровне одной строки создает ложную активную точку, чья интенсивность сравнима с интенсивностью настоящего сигнала. При высоком уровне шума это означает, что компоненте сигнала может соответствовать уже не первая, а одна из следующих компонент разложения. Оба дефекта ослабевают при уменьшении уровня шума: растекание вокруг правильного положения уменьшается, а смещение максимума может исчезать полностью.

На рис.~\ref{fig:row-row-components} показаны четыре характерных случая для одной и той же схемы \texttt{CSSA row.row}: смещение максимума, та же строка при меньшем уровне шума, растекание вокруг правильного положения и та же строка с растеканием при меньшем уровне шума.

\begin{figure}[H]
	\centering
	\includegraphics[width=\textwidth]{row_row_component_examples.pdf}
	\caption{Первая компонента для отдельных строк при \texttt{CSSA row.row}: смещение максимума, та же строка при $\sigma=0.05$, растекание вокруг правильного положения и та же строка при меньшем уровне шума}
	\label{fig:row-row-components}
\end{figure}

В первом примере максимум первой компоненты смещается в неверную точку. Во втором примере показана та же строка, но при меньшем уровне шума $\sigma=0.05$: смещение исчезает, и максимум возвращается к истинному положению пикселя сигнала. В третьем примере максимум остается правильным, однако вокруг него возникают дополнительные положительные значения меньшей величины. В четвертом примере показана та же строка, что и в третьем, но при меньшем уровне шума $\sigma=0.05$: растекание становится существенно слабее. Следовательно, оба указанных дефекта существенно зависят от уровня шумовой компоненты. Стоит заметить, что вместо уменьшения шума, что не применимо в реальных задачах, можно воспользоваться отсечением по фиксированному порогу или по порогу, задаваемому 0.9-квантилем. Последний вариант, хотя и может показывать хорошие результаты на синтетических данных, на практике оказывается ограниченно применим. Например, если одна прямая образована пикселями большой интенсивности, а другая пикселями малой интенсивности, то пороговая фильтрация может сохранить только первую прямую.

Растекание можно проанализировать и на уровне самой выделенной компоненты. На рис.~\ref{fig:smear-decomposition} показан пример строки, в которой максимум реконструкции остается в правильной позиции, но вокруг него возникают дополнительные положительные и отрицательные значения. Для идеального одиночного пикселя сигнала после дискретного преобразования Фурье модуль спектра должен быть постоянным по частотному индексу. В рассматриваемом примере амплитуда первой выделенной компоненты уже не является постоянной: она заметно меняется по частотам, убывает в центральной части и возрастает у краев. Следовательно, первая компонента уже не соответствует чистой комплексной экспоненте, и после обратного преобразования это приводит к растеканию сигнала по соседним пикселям. Фаза в данном случае не дает столь же простой интерпретации, однако также показывает отклонение от идеальной однокомпонентной структуры.

\begin{figure}[H]
	\centering
	\includegraphics[width=0.88\textwidth]{smear.pdf}
	\caption{Растекание в одной строке: реконструкция первой компоненты, ее амплитуда и фаза в спектральной области, а также вещественная и мнимая части вектора $U[,1]$}
	\label{fig:smear-decomposition}
\end{figure}

\section{Смещение максимума и неправильный выбор компоненты}

Для понимания природы смещения максимума рассмотрим однострочный эксперимент с фиксированной реализацией шума. Для одного и того же шумового вектора последовательно перебирались положения пикселя сигнала, пока не был найден первый индекс, в котором максимум реконструкции по первой компоненте ушел с истинной позиции. Этот случай показан на рис.~\ref{fig:one-row-bad-case}.

\begin{figure}[H]
	\centering
	\includegraphics[width=0.82\textwidth]{one_row_bad_case_first_component.pdf}
	\caption{Плохой случай для одной строки: слева реконструкция сигнала с максимумом в неверной точке, справа реконструкция того же шума без сигнала}
	\label{fig:one-row-bad-case}
\end{figure}

В найденной реализации пиксель сигнала расположен в позиции $38$, однако максимум реконструкции по первой компоненте оказывается в позиции $100$. При этом для того же фиксированного шума без сигнала максимум также достигается в позиции $100$. Следовательно, в данном случае первая компонента описывает прежде всего ряд, порожденный конкретной реализацией белого шума, а не полезным сигналом.

Вектор $U$ для настоящего сигнала был третьим по убыванию сингулярных значений, поэтому в реконструкцию по первой компоненте он не попал. Рассмотрим реконструкцию по третьей компоненте и соответствующие вещественную и мнимую части:

\begin{figure}[!htbp]
	\centering
	\includegraphics[width=0.82\textwidth]{one_row_bad_case_third_component.pdf}
	\caption{Реконструкция по третьей компоненте и вещественная с мнимой частями вектора $U[,3]$ для случаев \texttt{bad case} и \texttt{noise only}}
	\label{fig:one-row-bad-case-vectors}
\end{figure}

На рис.~\ref{fig:one-row-bad-case-vectors} видно, что при реконструкции по третьей компоненте максимум в случае \texttt{bad case} возвращается в правильную позицию. В то же время сам вектор $U[,3]$ имеет более регулярную колебательную структуру, чем первая компонента в плохом случае.

Если выполнить реконструкцию не по первой, а по третьей компоненте, максимум возвращается в позицию $38$. Следовательно, смещение максимума возникает не из-за малого числа компонент в реконструкции, а из-за неправильного выбора самой компоненты.

\section{От построчной реконструкции к входу преобразования Хафа}

Для метода Хафа несущественны сами значения интенсивности каждого пикселя, а множество точек, которые будут участвовать в голосовании. С учетом введенного в главе~1 нижнего порога это множество имеет вид
\[
E_{\varepsilon}(\widehat{\mathbf{S}})
=
\{(i,j)\colon \widehat{S}_{ij}>\varepsilon\}.
\]
Поэтому вместе с изображением после обработки \texttt{CSSA row.row} необходимо смотреть пороговый бинаризованный вариант, который фактически поступает на вход преобразования Хафа.

На рис.~\ref{fig:row-row-large-case} показана одна из проблемных реализаций, в которой после \texttt{CSSA row.row} сохраняются слабые горизонтальные следы.

\begin{figure}[!htbp]
	\centering
	\includegraphics[width=0.84\textwidth]{row_row_problem_case.pdf}
	\caption{Проблемная реализация: после \texttt{CSSA row.row} помимо основной прямой сохраняются слабые горизонтальные следы}
	\label{fig:row-row-large-case}
\end{figure}

После перехода к множеству активных пикселей именно эти слабые следы становятся важными. Бинаризованный вход для преобразования Хафа при малом пороге изображен на рис.~\ref{fig:row-row-binary-hough}.

\begin{figure}[!htbp]
	\centering
	\includegraphics[width=0.6\textwidth]{row_row_binary_hough_input.pdf}
	\caption{Бинаризованный результат \texttt{CSSA row.row}, то есть фактический вход преобразования Хафа}
	\label{fig:row-row-binary-hough}
\end{figure}

Таким образом, во время преобразования Хафа в аккумуляторной матрице будет накапливаться информация о потенциальной горизонтальной линии, поскольку бинаризованные остатки в отдельных строках образуют горизонтальные следы, что может приводить к ошибке оценивания.

\section{Статистический выбор порога активного множества}

Предыдущий пример показывает, что порог $\varepsilon$ нельзя выбирать только визуально. На восстановленном изображении слабые остатки могут казаться несущественными, однако после бинаризации каждый такой пиксель получает тот же статус активной точки, что и пиксель настоящей прямой. Поэтому порог должен быть связан с уровнем остаточного шума после предобработки.

Для одной строки после ДПФ можно записать
\[
\mathbf{Z}
=
F\mathbf{X}
=
\mathbf{S}+\boldsymbol{\eta},
\qquad
\boldsymbol{\eta}=F\boldsymbol{\varepsilon},
\qquad
\operatorname{Cov}(\boldsymbol{\eta})=N\sigma^2 I_N.
\]
Пусть после применения \texttt{CSSA row.row} получена оценка сигнальной части $\widehat{\mathbf{S}}$ в частотной области. Тогда остаток
\[
\mathbf{r}=\mathbf{Z}-\widehat{\mathbf{S}}
\]
можно рассматривать как приближение шумовой части в спектральной области. Если в остатке действительно доминирует шум, то его можно перевести обратно в исходную область:
\[
\widehat{\boldsymbol{\varepsilon}}
=
F^{-1}\mathbf{r}
=
\frac{1}{N}F^{*}\mathbf{r}.
\]
В случае чисто шумового остатка ковариация восстановленного шума равна
\[
\operatorname{Cov}(\widehat{\boldsymbol{\varepsilon}})
=
\frac{1}{N^2}F^{*}(N\sigma^2 I_N)F
=
\sigma^2 I_N,
\]
поскольку $F^{*}F=N I_N$. Значит, после обратного преобразования масштаб шума возвращается к исходному масштабу, и по значениям $\widehat{\boldsymbol{\varepsilon}}$ можно оценивать уровень остаточных случайных колебаний.

Практически это приводит к следующему правилу. Для каждой строки вычисляется остаток после CSSA-реконструкции, затем этот остаток переводится обратным ДПФ в исходную область. По полученным значениям оценивается стандартное отклонение остаточного шума, например
\[
\widehat{\sigma}_{\mathrm{res}}
=
\operatorname{sd}(\widehat{\boldsymbol{\varepsilon}}).
\]
После этого порог можно выбрать как высокий квантиль распределения модуля остаточного шума:
\[
\tau
=
q_{1-\alpha}\bigl(|\widehat{\boldsymbol{\varepsilon}}|\bigr),
\]
или, в более простой параметрической записи, как $\tau=c_{\alpha}\widehat{\sigma}_{\mathrm{res}}$, где $c_{\alpha}$ задает требуемую вероятность отсечения шумовых значений. Далее перед преобразованием Хафа используется активное множество
\[
E_{\tau}(\widehat{\mathbf{S}})
=
\{(i,j)\colon \widehat{S}_{ij}>\tau\}.
\]

Такой выбор порога отделяет этап подавления шума от этапа голосования в пространстве Хафа. CSSA отвечает за восстановление структуры прямой, а дополнительный порог удаляет слабые положительные остатки, которые не несут устойчивой геометрической информации, но способны создавать ложные максимумы в аккумуляторной матрице. Фиксированный порог, например $\tau=0.1$, можно использовать как простую экспериментальную версию этого шага, однако статистически мотивированный порог предпочтительнее, поскольку адаптируется к уровню остаточного шума.
